\section{Justificación de la propuesta}

\noindent\textbf{Interés académico, científico e industrial.} La propuesta atiende una necesidad verificable del Laboratorio de Mecatrónica de EAFIT y de talleres de prototipado, cuyas alternativas comerciales son manuales y no repetibles, o sobredimensionadas y costosas \cite{shannon_line,mic_auto}. Automatizar el proceso lo vuelve reproducible y trazable, reduce el desperdicio y libera tiempo calificado. Científicamente, exige cuantificar la relación entre la historia térmica local y la recuperación elástica del PMMA, problema de física de polímeros con tratamiento constitutivo no trivial \cite{gilormini2010} e impacto directo sobre la precisión alcanzable.

\noindent\textbf{Áreas de conocimiento involucradas.} Transferencia de calor \cite{incropera2011}; física de polímeros \cite{osswald2012,gilormini2010,yan2023}; mecánica de láminas delgadas; instrumentación y metrología \cite{iso17025,gum2008}; y control y sistemas embebidos \cite{astrom2006}.
